\section{Intellectuality and Spirituality Redux}

\begin{quotex}
The ``true" and the ``desired" must find their synthesis in the ``beautiful", for it is only in the beautiful that the urge to play renders the burden of the ``true" or the ``just" light and raises at the same time the darkness of instinctive forces to the level of light and consciousness. In other words, he who sees the beauty of that which he recognizes as true cannot fail to love it — and in loving it the element of constraint in the duty prescribed by the true will disappear: duty becomes a delight. \flright{\textsc{Valentin Tomberg}}

\end{quotex}
Since we began this sequence of gnosis meetings with the ideal of the alchemical marriage of intellectuality and spirituality several months ago, and now that we are finished for this year, it is worthwhile to circle back around on the topic.

\paragraph{Meditation}
Years ago, I learned centering prayer from Fr. Thomas Keating, not personally, but from some cassette tapes. In this brief video\footnote{\url{https://www.youtube.com/embed/pkDFaLdRck0}}, he asserts the necessity for daily meditation. For our purposes, note particularly his explanation that engaging with spiritual friends is an adequate substitute for a spiritual director. That — if you haven't figured it out that by now — is why we choose to work in groups. Hence, regular attendance is important, not just for yourself, but also for the commitments you've made to others.

Meditation, just like riding a bike or swimming, cannot be explained intellectually. One must simply begin. However, once the practice is established, you can get feedback from your spiritual companions. Just as your biking or swimming can improve, so can your meditating. After all, it is the first step to becoming a Bodhisattva.

\paragraph{The Hermetist and the Hermetic Path}
The Hermetic path is a gift, if you are called to it. It requires intelligence, resources, and time, so it is certainly not an option for everyone. However — and this is important — it is not ``superior" to, or an alternative to, the exoteric path, it is simply our path. It would be a mistake to force these views on exoteric practitioners, or to use them as debating points. The exoteric path is perfectly adequate for salvation and a life of sanctity. There is a reason Hermetic groups used to be secret and closed to outsiders.

The Hermetist often used to masquerade as a trader or street performer; the latter is the primary meaning of the first Tarot card, Le Bateleur. Under such cover, they could travel from town to town, allowing them to meet with local groups without attracting attention. Since books were heavy and expensive, the teachings were conveyed in diagrams (as in \emph{Gnosis}), or even a deck of Tarot cards. On the one hand, they were compact and portable, but on the other, they required an accompanying oral teaching in order to be fully understood.

Years ago, I used to follow the books of \textbf{Carlos Castaneda}. I'm sure the teachings of the shaman Don Juan are still embedded in my soul somewhere, for better or for worse. The first thing to note is that it is difficult to find Don Juan. In one scene that I recall, Don Juan appeared at some government office on official business in a suit, just blending in with everyone else. This struck Castaneda, who had never seen him in that context, as something remarkable. Nowadays, shamans seem to be everywhere, peddling their books, courses, and so on. Those converts to some sort of spiritual life often feel they have to alter their outer appearance to be convincing. This is quite unlike Don Juan or even a Hermetist.

The real shaman Don Juan was indistinguishable from his surroundings. Therefore, you could not pick out the Hermetist in a room. He would look like everyone else and talk about the weather or sports to you. Only if you expressed some sort of interest in something deeper, might he open up to you.

St. Augustine tells us that God gives us everything we need for salvation. So if you are still searching, you may be missing the obvious. This might be as good as it gets for you.

\paragraph{The Quest for Novelty}
Some birds allegedly become fascinated by shiny objects on the ground, thereby forgetting the bugs that constitute their true nourishment. Similarly, we often become distracted by convoluted intellectual schemes or elaborate tableaux. Stay rooted in the core principles, and be clear about the difference between an allegory or secret, and a genuine mystery. Tomberg warns us about this temptation:

\begin{quotex}
Let us therefore not commit the error of wanting to ``explain" a symbol by reducing it to a few general abstract ideas. Let us also avoid the error of wanting to ``concretise" an abstract idea by clothing it in the form of an allegory. 

\end{quotex}
The Arcana are not allegories in which a card is said to ``represent" some qualities along the lines, say, of a work like John Bunyan's Pilgrim's Progress. In Tomberg's words:

\begin{quotex}
The Major Arcana of the Tarot are neither allegories nor secrets, because allegories are, in fact, only figurative representations of abstract notions, and secrets are only facts, procedures, practices, or whatever doctrines that one keeps to oneself for a personal motive, since they are able to be understood and put into practice by others to whom one does not want to reveal them. The Major Arcana of the Tarot are authentic symbols. They conceal and reveal their sense at one and the same time according to the depth of meditation. 

\end{quotex}
In other words, if the meaning seems to ``jump out" at you immediately, it may not be an ``authentic symbol" in this sense. Ultimately, we will come to understand the mystery, but only plunging into its depths. An allegory, on the other hand, tries to ``solve" the mystery:

\begin{quotex}
Just as the \emph{arcanum }is superior to the \emph{secret, }so is the \emph{mystery }superior to the Arcanum. The mystery is more than a stimulating ``ferment". It is a spiritual \emph{event }comparable to physical birth or death. It is a change of the entire spiritual and psychic motivation, or a complete change of the plane of consciousness. 

\end{quotex}
This cannot be overemphasized. A mystery cannot be resolved intellectually. Unless an arcane teaching leads to, or elicits, or acts as the midwife to, a spiritual transformation, it has not been understood. And worse, it may even become a distraction, with no more transformative power than a parlour game.

\paragraph{Cartesian Meditation}
Cartesian meditation, which is the search for clear and distinct ideas, is an intellectual task. \textbf{Rene Descartes} was a mathematical prodigy. As a boy, I was quite proficient in mathematics, although not at that level, so I enjoyed reading about the lives of the great mathematicians. Now, because of his intellect, the young Rene Descartes was pampered. Instead of being shooed out of bed in the morning, he was allowed to lounge, giving him the leisure to think.

There is no doubt that there can be great joy in the experience of intellectual insight or learning. Just watch the expression on a baby's face when he takes his first step and young children when they learn a new skill. I've seen people show great excitement in solving a puzzle or answering a question while watching some game show on TV. Adults still do crosswords or Sudoku just for the pleasure of it.

That feeling is magnified with more complex intellectual attainments, particularly in physics, maths, and metaphysics. For example, Newton's discovery of the equivalence of inertia and gravity is mind bending when it dawns on you, as is Descartes' discovery of the transformability of algebra and geometry. I personally can attest to the pleasures in physics and maths. One can struggle with an obscure maths problem, but, in a sudden insight, its solution simply appears, perhaps analogous to the experience of \emph{Yesod}. Of course, the study of metaphysics can lead to a sort of bliss, especially with the realization that certain ideas bring you oh so close to the very nature of God. At this point, the search for Truth becomes the delight in Beauty.

So, back to the young Rene: following his example. I will often lie in bed pondering some issue. Of course, Cartesian meditation is not the source. Rather, the real meditation reaches in the depths, often murky depths, not for the clarity of the atmosphere. Nevertheless, clear ideas are floating in the darkness of those depths, and they need to be coaxed out. Obviously, the discursive mind is required in order to turn those vague intuitions into text. That is the purpose of Cartesian meditation. Ultimately, however, there is not a shortage of ideas, but rather its opposite. There is actually an abundance of ideas, so cutting and pruning is necessary. Much more is discarded than is ever published.

If you allow your intellectual life to be nourished by the real nutriments hiding in the darkness, you will no longer be satisfied with dazzling baubles, word puzzles, or intellectual trivia. The goal of the intellectual life is to be a Sage, so seek the higher things like virtue, the life of reason, aesthetic beauty, the path of salvation, and the attributes of God.

\paragraph{Living in the Light of Tabor}
\begin{quotex}
Hermetism is an \emph{athanor }(``alchemical furnace") erected in the individual human consciousness, where the mercury of intellectuality undergoes transmutation into the gold of spirituality. St. Augustine acted as a Hermetist in transmuting Platonism into Christian thought. Similarly, St. Thomas Aquinas acted as a Hermetist in doing the same thing with Aristotelianism. Both of them accomplished the sacrament of baptism with respect to Greece's intellectual heritage. \flright{\textsc{Valentin Tomberg}, \emph{Letter on Justice}}

\end{quotex}
If the goal of the intellectual life is to become a Sage, the goal of the spiritual life is to become a Saint. Of course, the latter quest is foolishness to the Intellect. Hence, only the Fool can show us the way. Tomberg explains:

\begin{quotex}
The Arcanum ``The Fool" teaches the ``know-how" of passing from intellectuality, moved by the desire for knowledge, to the higher knowledge due to love. 

\end{quotex}
The Fool is also a Trickster, since there are two ways of sacrificing the intellect:

\begin{itemize}
\item It can submit itself to the service of Transcendental consciousness 
\item It can simply be abandoned 
\end{itemize}
Now, the temptation to simply abandon the intellect is quite strong. Deep meditation may be accompanied with intense sensations of pleasure, or even siddhis. The spiritual quest may then devolve to a quest for the repetition of such feelings as ends in themselves. Common practices such as postures, breath control, dancing, chanting, and so on may help focus the mind. When they fail, some schools resort to stronger practices such as drugs, alcohol, or extreme sexuality. However, the Spirit cannot be coerced by any sort of technique or mechanical practices.

Tomberg points to the Whirling Dervishes and Zen monks as those who have abandoned the intellect entirely. Some of this lies behind the Hesychast controversy. What concerned Barlaam was the ignorance and credulity of some of the monks, so, in compensation, he overemphasised the side of the Intellect. The monks, on the other hand, pointed out that the first disciples were simple men, not advanced scholars. Now that may be true in the Synoptic Gospels, but John's Gospel explicitly identifies Christ with the Logos behind the creation of the world. We take the middle path between Barlaam and Palamas.

\paragraph{The Work}
In our time apart, we could focus on all the themes of the past few years. We can be Holy Fools, yet still be intellectually competent. Our meditations should be on the life of Christ or something analogous; that is, something that requires an Active Imagination, not the passive imagination of a dream-like state. We concentrate without effort and have mastery over what thoughts and emotions are allowed to take hold in our consciousness.

To achieve the fusion of intellectuality and spirituality, we need to return up the Middle Pillar. That begins with the recognition of one's True Will and ends with the awareness of one's Real I. That is the gift of Integrity that was lost in the Fall.

\flrightit{Posted on 2017-05-26 by Cologero }
